\documentclass[10pt,a4paper]{amsart}
\usepackage[margin=19mm]{geometry}
\usepackage{amsmath,amssymb,mathtools}
\usepackage[scr=rsfs,cal=euler]{mathalfa}
\usepackage{xfrac}
\usepackage[unicode,hidelinks,bookmarks=false]{hyperref}
\newcommand{\ie}{i.e.\ }
\newcommand{\cf}{cf.\ }
\newcommand{\resp}{respectively\ }
\newcommand{\Subsection}{\subsection}
\providecommand{\nobreakdash}{\nobreak\hbox{-}}
\setlength{\emergencystretch}{2em}
\multlinegap0pt
\usepackage{bm}
\usepackage{bbm}
\usepackage{dsfont}
\usepackage{xspace}
\usepackage{cancel}
\usepackage{mathtools}
\usepackage{amsthm}
\makeatletter
\@ifundefined{Theorem}{\newtheorem{Theorem}{Theorem}}{}
\@ifundefined{Lemma}{\newtheorem{Lemma}[Theorem]{Lemma}}{}
\makeatother
\title[Topological and Metric Pressure for Singular Flows]{Topological and Metric Pressure for Singular Flows}
\author{Meijie Zhao, Xiao Wen}
\date{Source: arXiv:2510.27216v1 (CC BY 4.0)}
\begin{document}
\maketitle

\noindent\emph{Source and licence.} Topological and Metric Pressure for Singular Flows. Version arXiv:2510.27216v1 (CC BY 4.0), distributed under the Creative Commons Attribution 4.0 International licence, as recorded in the arXiv licence metadata for this exact version and shown on the abstract page https://arxiv.org/abs/2510.27216v1. Full rights notice: licenses/arxiv-2510.27216-CC-BY-4.0.txt.\par\smallskip

\noindent\textit{Editorial scope.} Counted unit: the Lemma whose source label is \texttt{lemma4.6} in the article (source0.tex lines 1063--1068, source label \texttt{lemma4.6}), delivered together with the article's own complete original proof of it (source0.tex lines 1069--1261, one \texttt{proof} environment of 193 lines). No mathematics is written, repaired or re-derived: statement and proof are the article's own text line for line. Required context retained uncounted: none. Every definition, display and label of those lines is retained unchanged. Omitted from this package and not redistributed: 1--1062, 1262--1872. Nothing inside a retained excerpt is omitted or altered: every retained body is delivered exactly as the article's lines are written, so no omission receipt of any kind accompanies this package. Citations: the retained excerpts carry no citation command at all, so no bibliography entry of the article is transcribed and no reference is redistributed. Reference adaptation: none was needed and none was made; every reference inside the delivered unit resolves inside it, so no unresolved reference or citation is produced. Printed slips kept literal: no printed word, symbol, hypothesis or number of the counted statement or of its proof was altered. Layout and class: the article's own class is \texttt{article} with packages including hyperref, caption, graphicx, latexsym, amssymb, amsthm, indentfirst, amsmath, color, xcolor, geometry; the wrapper is the lane's pinned \texttt{amsart} editorial wrapper at 10pt on A4, which restates the theorem-like environments the retained text needs and adds the article's own macro definitions that the retained text needs (none), with extra wrapper packages (bm, bbm, dsfont, xspace, cancel, mathtools). The delivered numbering is the unit's own. Assets: no retained span contains a figure environment, a graphics-inclusion command, a PDF-page inclusion, a table or a TikZ picture; no image file is copied into this package, the package declares no asset dependency and no image file is redistributed with it. Licence: version arXiv:2510.27216v1 of this article is distributed under the Creative Commons Attribution 4.0 International licence, as recorded in the arXiv licence metadata for this exact version and shown on the abstract page https://arxiv.org/abs/2510.27216v1. A full rights notice is delivered at licenses/arxiv-2510.27216-CC-BY-4.0.txt. Characters: every retained excerpt is pure ASCII, so the delivered engine, XeLaTeX with its default Latin Modern text font, reports no missing character.
\begin{Lemma}\label{lemma4.6}
Suppose that $\mu$ is an ergodic measure of \( \phi_\tau \) for some constant $\tau>0$. Let \( \xi = \{ C_1, C_2, \ldots, C_N \} \) (\( N \geq 3 \)) be a finite measurable partition of \( M \) with \( \mu(\partial \xi) = 0 \). Then for any \( \delta \in (0, 1) \) and \( f \in C(M, \mathbb{R}) \), we have
\[
\lim_{\gamma \to 0} \liminf_{n \to \infty} \frac{1}{n} \log N_1^\mu(\delta, n\tau, \gamma, f, X) \geq \frac{1}{\tau} h_\mu(\phi_\tau, \xi) + \int f \, {\rm d}\mu .
\]
\end{Lemma}

\begin{proof}
Let \( 0 < \varepsilon < \frac{1 - \delta}{2} \) be an arbitrarily given positive constant. Since \( \mu(\partial \xi) = 0 \), there exists a neighborhood \( U \) of \( \partial \xi \) such that \( \mu(U) < \frac{\varepsilon^2}{8} \). Denote by \( \chi_U \) the characteristic function of \( U \). For any \( n \in \mathbb{N} \), define
\[
B_{n, \varepsilon} = \left\{ x \in M : \frac{1}{n} \sum_{i=0}^{n-1} \chi_U(\phi_{i\tau}(x)) < \frac{\varepsilon }{2} \right\}.
\]
\noindent
\textbf{Claim 1. } \(\mu(B_{n, \varepsilon}) > 1 - \frac{\varepsilon}{4} \).

\bigskip


\noindent
\begin{proof} Note that \( \int_M \chi_U(x) \, {\rm d}\mu = \mu(U) < \frac{\varepsilon^2}{8} \). Since \( \mu \) is an invariant measure for \( \phi_\tau \), we have
\[
\int_M \chi_U ( \phi_{i\tau }(x)) \, {\rm d}\mu = \int_M \chi_U(x) \, {\rm d}\mu < \frac{\varepsilon^2}{8}, \quad \forall\, 0 \leq i \leq n-1.
\]
Thus
\[
\sum_{i=0}^{n-1} \int_M \chi_U( \phi_{i\tau }(x)) \, {\rm d}\mu < \frac{n }{8}\varepsilon^2.
\]
On the other hand
\begin{align*}
\sum_{i=0}^{n-1} \int_M \chi_U ( \phi_{i\tau }(x)) \, {\rm d}\mu &= \int_M \left( \sum_{i=0}^{n-1} \chi_U(\phi_{i\tau }(x)) \right) {\rm d}\mu\geq \int_{M \setminus B_{n, \varepsilon}} \left( \sum_{i=0}^{n-1} \chi_U(\phi_{i\tau }(x)) \right) {\rm d}\mu \\&\geq\int_{M \setminus B_{n, \varepsilon}} \frac{n\varepsilon}{2}{\rm d}\mu= \frac{n \varepsilon}{2} \mu(M \setminus B_{n, \varepsilon}).
\end{align*}
\noindent
Thus \( \frac{n\varepsilon}{2} \mu(M \setminus B_{n, \varepsilon}) < \frac{n}{8} \varepsilon^2 \), then \( \mu(M \setminus B_{n, \varepsilon}) < \frac{\varepsilon}{4} \).
\end{proof}

Also, for \( \mu \)-a.e. \( x \in M \)
\[
\lim_{n \to \infty} \frac{1}{n\tau} \int_0^{n\tau} f(\phi_{ t}(x)) \, {\rm d}t= \int f \, {\rm d}\mu.
\]

\noindent
By Egoroff's Theorem, there exists \( B \subseteq M \) such that \( \mu(B) > 1 - \frac{\varepsilon}{4} \) and \( \frac{1}{n\tau} \int_0^{n\tau} f(\phi_{ t}(x)) \, {\rm d}t \) converges uniformly to \( \int f \, {\rm d}\mu \) as \( n \to \infty \) on \( B \). That is, there exists \( N_1 \) such that for any \( n > N_1 \), we have
\[
\left| \frac{1}{n\tau} \int_0^{n\tau} f(\phi_{ t}(x)) \,{\rm d} t - \int f \, {\rm d}\mu \right| < \varepsilon
\]
for any \( x \in B \).

Let \( A_{n, r}^- = \left\{ A \in  \xi_n : \mu(A) < e^{-n(h_\mu(\phi_\tau, \xi) - \varepsilon)} \right\} \), where \( \xi_n = \xi \vee \phi_{\tau}^{-1} \xi \vee \cdots \vee \phi_{\tau}^{-(n-1)} \xi \). Similarly, by the Shannon-McMillan-Breiman Theorem, there exists \( N_2 \) such that for any \( n > N_2 \), we have
\[
\mu\left( \bigcup_{A \in A^-_{n, r}} A \right) > 1 - \frac{\varepsilon}{4}.
\]

Fix arbitrary $f\in C(M, \mathbb{R})$. We can take \( \gamma_0 > 0 \) such that if \( d(x, y) < \gamma_0 \|X(x)\| \), then the followings hold:
\begin{enumerate}
\item[1.]
    \(
    |f(x) - f(y)| < \varepsilon,
    \)
\item[2.]
    Either  \( \{x, y\} \subseteq U \) or  \( \{x, y\} \subseteq C_i \) for some \( i = 1, 2, \ldots, N \).
\end{enumerate}

The first item in above follows from the uniform continuity of \( f \). If there were no such \( \gamma_0 \) such that the second item holds, then we could take \( x_k, y_k \) containing in different \( C_i \) respectively, and \( x_k, y_k \) not both containing in \( U \), with \( d(x_k, y_k) \to 0 \). By the compactness of $M$, without loss of generality, assume that \( x_k, y_k \) converge (to the same point). Then their limit point is an element in \( \partial{\xi} \). When \( k \) is sufficiently large, \( x_k, y_k \) would both belong to \( U \), which is a contradiction.

Fix \( n > \max\{N_1, N_2\} \) , \( 0 < \gamma < \gamma_0 \). Assume that there exists a finite set \( F \subseteq M \setminus \text{Sing}(X) \) satisfying
\[
\mu\left( \bigcup_{x \in F} B_1^*(x, n\tau, \gamma, X) \right) > 1 - \delta.
\]
Then
\[
\mu\left( \left( \bigcup_{x \in F} B_1^*(x, n\tau, \gamma, X) \right) \cap \left( \bigcup_{A \in \mathcal{A}^-_{n, r}} A \right) \cap B_{n, \varepsilon} \cap B \right) > 1 - \delta - \frac{\varepsilon}{4} - \frac{\varepsilon}{4} - \frac{\varepsilon}{4}> \frac{1 - \delta}{4}.
\]

Let
\(
Q = \left( \bigcup\limits_{x \in F} B_1^*(x, n\tau, \gamma, X) \right) \cap \left( \bigcup\limits_{A \in A^-_{n, r}} A \right) \cap B_{n, \varepsilon} \cap B \), and set
\[
\mathcal{E} = \left\{ A \in \xi_{n} : A \cap Q \neq \emptyset \right\}.
\]
Then \( \mathcal{E} = \left\{ A \in A^-_{n,r} : A \cap \left( \bigcup\limits_{x \in F} B_1^*(x, n\tau, \gamma, X) \right) \cap B_{n,\varepsilon} \cap B \neq \emptyset \right\} \). For \( A \in \mathcal{E} \), since \( A \in {A}_{n,r}^- \), we have \( \mu(A) < e^{-n(h_\mu(\phi_\tau, \xi) - \varepsilon)} \), while \( \mu\left( \bigcup\limits_{A \in \mathcal{E}} A \right) > \mu(Q) > \frac{1 - \delta}{4} \). Thus we have
\[
\#\mathcal{E} > \frac{1 - \delta}{4} e^{n(h_\mu(\phi_\tau, \xi) - \varepsilon)}.
\]

For any \( x \in M \), there exists a unique map \( \omega_x: \{0, 1, \ldots, n - 1\} \to \{1, 2, \ldots, N\} \) such that \( \phi_{i\tau}(x)\in C_{\omega_x(i)} \). It is easy to see that if
\[
y \in \xi_n(x) = \xi(x) \cap \phi_{-\tau}(\xi(\phi_\tau(x))) \cap \cdots \cap \phi_{-(n-1)\tau}(\xi(\phi_{(n-1)\tau}(x))),
\]
then we have \( \omega_x = \omega_y \), where \( \xi(x) \) is the element in \( \xi \) containing \( x \).

Let \( F' = \left\{ x \in F : B_1^*(x, n\tau, \gamma, X) \cap B_{n,\varepsilon} \cap B \neq \emptyset \right\} \). For any \( x \in F' \), let
\[
\mathcal{E}_x = \left\{ A \in {A}_{n,r}^- : A \cap B_1^*(x, n\tau, \gamma, X) \cap B_{n,\varepsilon} \cap B \neq \emptyset \right\}. \tag{4}
\]
Then \( \mathcal{E} \subseteq \bigcup\limits_{x \in F'} \mathcal{E}_x \).

We have the following estimation for the number of elements in \( \mathcal{E}_x \).

\bigskip
\noindent
\(\textbf{Claim 2. }\) {\it For any \( x \in F' \), we have \( \#\mathcal{E}_x \leq P(N, n, \frac{\varepsilon}{2}) \).}



\begin{proof}
Let \( A \in \mathcal{E}_x \). Note that for any \( y, z \in A \), we  have \( \omega_y = \omega_z \), thus we can denote this common sequence as \( \omega_A \). Since \( A \in \mathcal{E}_x \), there exists \( y \in A \cap B_1^*(x, n\tau, \gamma, X) \cap B_{n, \varepsilon} \cap B \), and \( \omega_y = \omega_A \).

Since \( y \in B_1^*(x, n\tau, \gamma, X) \), then for all \( 0 \leq i \leq n - 1 \), we have
\[
d(\phi_{i\tau}(x), \phi_{i\tau}(y)) < \gamma \| X(\phi_{i\tau}(x)) \| < \gamma_0 \| X(\phi_{i\tau}(x)) \|.
\]
Thus, \( \phi_{i\tau}(x) \) and \( \phi_{i\tau}(y)\, (0\leq i\leq n-1)\) either both belong to the same element of \(\xi\) or both belong to \( U \). That is, if \( \omega_x(i) \neq \omega_y(i) \), then \( \chi_U(\phi_{i\tau}(x)) = \chi_U(\phi_{i\tau}(y)) = 1 \). Thus we have
\[
1 - \delta_{\omega_x(i), \omega_y(i)} \leq \chi_U(\phi_{i\tau}(x))
\]
for any $0\leq i\leq n-1$. Then
\[
\rho_{N,n}(\omega_x, \omega_y) = \frac{1}{n} \sum_{i = 0}^{n - 1} (1 - \delta_{\omega_x(i), \omega_y(i)}) \leq \frac{1}{n} \sum_{i = 0}^{n - 1} \chi_U(\phi_{i\tau}(x)) < \frac{\varepsilon}{2}.
\]
This proves that \( \rho_{N,n}(\omega_x, \omega_A) < \frac{\varepsilon}{2} \), and then \( \omega_A \in B\left( \omega_x, \frac{\varepsilon}{2}; N, n \right) \).

Therefore, when \(A\in \mathcal{E}_x \), we always have \( \omega_A \in B\left( \omega_x, \frac{\varepsilon}{2}; N, n \right) \). Note that for any $A, B\in \xi_n$, \( A=B \) if and only if \( \omega_A = \omega_B \) . We can see that \( \#\mathcal{E}_x \leq \# B\left( \omega_x, \frac{\varepsilon}{2}; N, n \right) = P\left( N, n, \frac{\varepsilon}{2} \right) \). This finishes the proof of the claim.
\end{proof}

By the claim, we have
\[
\frac{1 - \delta}{4} e^{n(h_\mu(\phi_\tau, \xi) - \varepsilon)} < \#\mathcal{E} = \#\left( \bigcup_{x \in F'} \mathcal{E}_x \right) \leq \# F' \cdot P\left( N, n, \frac{\varepsilon}{2} \right).
\]
Hence we have
\[
\# F' \geq \frac{1 - \delta}{4} e^{n(h_\mu(\phi_\tau, \xi) - \varepsilon)} \cdot P\left( N, n, \frac{\varepsilon}{2} \right)^{-1}.
\]

For any \( x \in F' \), since \( B_1^*(x, n\tau, \gamma, X) \cap B_{n, \varepsilon} \cap B \neq \emptyset \), there exists
\(
y \in B_1^*(x, n\tau, \gamma, X) \cap B_{n, \varepsilon} \cap B.
\)
For any $t\in[0, n\tau]$, because
\[
d(\phi_t(x), \phi_t(y)) < \gamma\| X(\phi_t(x)) \| < \gamma_0 \| X(\phi_t(x)) \|,
\]
we have \(
|f(\phi_t(x)) - f(\phi_t(y))| < \varepsilon.
\)

Note that for any \( y \in B \) and any \( n > N_1 \), we have
\[
\left| \frac{1}{n\tau} \int_0^{n\tau} f(\phi_t(y)) {\rm d}t - \int f {\rm d}\mu \right| < \varepsilon.
\]
Thus we have
\[
\begin{aligned}
&\left| \frac{1}{n\tau} \int_0^{n\tau} f(\phi_t(x)) {\rm d}t - \int f {\rm d}\mu \right| \\
\leq& \frac{1}{n\tau} \left( \left| \int_0^{n\tau} f(\phi_t(x)) {\rm d}t - \int_0^{n\tau} f(\phi_t(y)) {\rm d}t \right| + \left| \int_0^{n\tau} f(\phi_t(y)) {\rm d}t-  n\tau\int f {\rm d}\mu \right| \right)\\
<& 2\varepsilon.
\end{aligned}
\]
And then
\[
\int_0^{n\tau} f(\phi_t(x)) {\rm d}t> n\tau \left( \int f {\rm d}\mu - 2\varepsilon \right).
\]
This follows that
\[
\begin{aligned}
\sum_{x \in F} e^{\int_0^{n\tau} f(\phi_t(x)) {\rm d}t} &\geq \sum_{x \in F'} e^{\int_0^{n\tau} f(\phi_t(x)) {\rm d}t} \geq \sum_{x \in F'} e^{n\tau \left( \int f {\rm d}\mu - 2\varepsilon \right)} \geq \# F' \cdot e^{n\tau \left( \int f {\rm d}\mu - 2\varepsilon \right)} \\
&\geq \frac{1 - \delta}{4} e^{n(h_\mu(\phi_\tau, \xi) - \varepsilon)} \cdot P\left( N, n, \frac{\varepsilon}{2} \right)^{-1} \cdot e^{n\tau \left( \int f {\rm d}\mu - 2\varepsilon \right)}.
\end{aligned}
\]

This proves that if \( n > \max\{N_1, N_2\}, \,0<\gamma<\gamma_0 \), then we have
\[
N_1^\mu(\delta, n\tau, \gamma, X, f) \geq \frac{1 - \delta}{4} e^{n(h_\mu(\phi_\tau, \xi) - \varepsilon)} P\left( N, n, \frac{\varepsilon}{2} \right)^{-1} e^{n\tau \left( \int f {\rm d}\mu - 2\varepsilon \right)}.
\]

\noindent
It follows that
\[
\begin{aligned}
&\liminf_{n \to \infty} \frac{1}{n\tau} \log N_1^\mu(\delta, n\tau, \gamma, X, f) \\
\geq& \liminf_{n \to \infty} \frac{1}{n\tau} \left( n(h_\mu(\phi_\tau, \xi) - \varepsilon) + n\tau \left( \int f {\rm d}\mu -2\varepsilon \right)+\log(1 - \delta) - \log 4 - \log P\left( N, n, \frac{\varepsilon}{2} \right) \right) \\
=& \frac{1}{\tau} h_\mu(\phi_\tau, \xi) + \int f {\rm d}\mu - \frac{\varepsilon}{\tau} - 2\varepsilon -\frac{1}{\tau} \left(\frac{\varepsilon}{2} \log(N - 1) - \frac{\varepsilon}{2} \log \frac{\varepsilon}{2} - \left( 1 - \frac{\varepsilon}{2} \right) \log \left( 1 - \frac{\varepsilon}{2} \right)\right).
\end{aligned}
\]

\noindent
Thus
\[
\begin{aligned}
\lim_{\gamma\to 0} \liminf_{n \to \infty} \frac{1}{n\tau} \log N_1^\mu(\delta, n\tau, \gamma, X, f) &\geq \frac{1}{\tau} h_\mu(\phi_\tau, \xi) + \int f {\rm d}\mu - \frac{\varepsilon}{\tau} - 2\varepsilon -\frac{\varepsilon}{2\tau} \log(N - 1) \\
&+ \frac{\varepsilon}{2\tau} \log \frac{\varepsilon}{2} +\frac{1}{\tau} \left( 1 - \frac{\varepsilon}{2} \right) \log \left( 1 - \frac{\varepsilon}{2} \right).
\end{aligned}
\]

\noindent
By letting \( \varepsilon \to 0 \), we can get
\[
\lim_{\gamma\to 0} \liminf_{n \to \infty} \frac{1}{n\tau} \log N_1^\mu(\delta, n\tau, \gamma, X, f) \geq \frac{1}{\tau} h_\mu(\phi_\tau, \xi) + \int f {\rm d}\mu.
\]
This finishes the proof of the lemma.
\end{proof}

\end{document}
